% =====================================================================
% paper_skeleton.tex  —  MDPI MTI section scaffold
% =====================================================================
% HOW TO USE: this is NOT a standalone file. Start from the official MDPI
% Overleaf template (Definitions/mdpi.cls), then paste these \section blocks
% into the template body. Keep the template's front-matter (\Title, \Author,
% \abstract, \keyword) and back-matter (\authorcontributions, \funding, etc.).
% Citation keys below match references.bib. Replace every <NOTE ...> with prose
% (ask Claude to draft each one — see ../README.md Part C).
% =====================================================================

% ---- FRONT MATTER (fill in the template's macros, shown here for reference) ----
% \Title{Leakage-Free Pedestrian Crossing-Intention Prediction from Bounding
%        Boxes and Ego-Speed: A Parsimonious BiLSTM Baseline on PIE}
% \Author{...}  \address{...}  \corres{...}

% \abstract{<NOTE ~200 words: (1) safety problem + intention prediction; (2) the
% temporal-leakage finding in common PIE protocols; (3) our leakage-free,
% crossing_point-anchored protocol; (4) a 2-stream (bbox + ego-speed) BiLSTM reaches
% AUC 0.932 (95% CI 0.92-0.95) -- top of the standard-protocol band -- at 0.575 ms per
% window; (5) shown with bootstrap CIs, LOSO, multi-seed ablations, a documented
% hyperparameter search, and detector-in-the-loop evaluation. Conclude: minimal
% modality is enough; the pipeline is detection-bound, not prediction-bound.>}

% \keyword{pedestrian crossing intention; intention prediction; PIE dataset; BiLSTM;
% temporal leakage; ADAS; ego-vehicle speed; real-time inference}

% =====================================================================
\section{Introduction}
% <NOTE: safety motivation -- pedestrian fatalities, ADAS need to anticipate crossing
% BEFORE it happens; intention != trajectory. Cite PIE \cite{rasouli2019pie}.>
% <NOTE: the gap. (a) Many PIE works never check TEMPORAL LEAKAGE -- windows anchored
% at track end can already contain crossing frames (Issue 1). (b) Recent SOTA is
% heavy/multimodal (pose, optical flow, depth, semseg). (c) Single-split, single-seed,
% GT-box evaluation is common. Cite PCPA \cite{kotseruba2021benchmark}, GTransPDM
% \cite{gao2024gtranspdm}, PIP-Net \cite{azarmi2025pipnet}.>
\paragraph{Contributions.}
\begin{itemize}
  \item % <NOTE: a temporal-leakage audit of the common PIE protocol + a leakage-free
        % re-extraction anchored at PIE's crossing_point (0% leakage verified). Issues 1-2.>
  \item % <NOTE: a parsimonious 2-stream (bbox + ego-speed) BiLSTM that matches top
        % standard-protocol AUC (0.932) at real-time latency (0.575 ms/window). Issues 3, 9.>
  \item % <NOTE: a rigorous evaluation -- bootstrap CIs, LOSO, multi-seed ablations,
        % a documented hyperparameter search, and detector-in-the-loop. Issues 4-8, 10.>
\end{itemize}
% <NOTE: 1-sentence roadmap of the paper.>

% =====================================================================
\section{Related Work}
% <NOTE: PIE crossing-prediction landscape. Build from
% ../journal_prep/issue3_baseline_comparison/03_baseline_comparison.md:
% PCPA \cite{kotseruba2021benchmark} (4-stream, AUC 0.86, defines the protocol);
% Pedestrian Graph+ \cite{cadena2022pedgraph}; PIT/IntFormer; BiPed/PedFormer
% \cite{rasouli2023pedformer}; GTransPDM \cite{gao2024gtranspdm} (pose+graph, 0.87);
% PIP-Net \cite{azarmi2025pipnet} (7 modalities, 0.897).>
% <NOTE: Occlusion-Aware Diffusion \cite{liu2025occlusion} as a PRECEDENT that
% bbox+ego-velocity is a legitimate minimal modality -- but NOT a comparison row
% (occluded-only, ~1-frame-ahead TTE). State the caveat.>
% <NOTE: the field-wide gaps we close -- from 04_positioning_vs_prior_work.md:
% leakage unchecked, single split, point estimates, single-seed ablations, GT-box
% assumption, unjustified hyperparameters.>

% =====================================================================
\section{Materials and Methods}

\subsection{The PIE Dataset}
% <NOTE: PIE \cite{rasouli2019pie} -- on-board video, set01-06, bbox + vehicle OBD
% ego-speed, crossing labels. Fixed split: train set01/02/04, val set05/06, test set03.>

\subsection{Temporal Leakage and the Leakage-Free Protocol}
% <NOTE: THE distinctive contribution. Common protocol anchors the window at the last
% annotated frame minus TTE; this can overlap frames where the pedestrian is ALREADY
% crossing -- the task collapses to "detect in-progress crossing." Issue 1: 67.9% of
% crossers leaked under the old anchor. Fix: anchor at PIE's own crossing_point
% attribute, slide obs_len=16 windows at 50% overlap with TTE in [30,60] frames before
% the crossing point -> 0% leakage verified (Issue 2). N grows 1,389 -> 4,906.>

\subsection{Features and Preprocessing}
% <NOTE: 5-D feature [x1,y1,x2,y2,vehicle_speed] in raw PIE pixel coords (1920x1080),
% train-only standardisation (mean/std from train split). obs_len=16, threshold 0.5.>

\subsection{Model Architecture}
% <NOTE: BiLSTMIntentPredictor -- input projection (Linear 5->64 + ReLU), 2-layer
% bidirectional LSTM (hidden 128), linear head on the last step. ~0.6M params. Capacity
% justified in Results (Issue 7: hidden + depth ablation).>

\subsection{Training and Hyperparameter Selection}
% <NOTE: Adam lr=1e-3, weight_decay=1e-5, batch 32, BCEWithLogits with pos_weight=1.682
% (clean train balance), early stop on val AUC (patience 15), best-val checkpoint, test
% touched once. Reproducibility: seed; 5 seeds for headline. Issue 8: a documented
% 36-config grid search (val-only selection) CONFIRMS this hand-set config -> not arbitrary.>

\subsection{Live Pipeline (Perception to Prediction)}
% <NOTE: YOLO26-M \cite{yolo} person detection + ByteTrack \cite{zhang2022bytetrack}
% multi-object tracking feed 16-frame bbox windows + ego-speed into the BiLSTM. Used
% for the latency (Issue 9) and detector-in-the-loop (Issue 10) experiments.>

% =====================================================================
\section{Results}

\subsection{Main Result and Confidence Intervals}
% <NOTE: AUC 0.932 +/- 0.011 (5 seeds), 95% bootstrap CI [0.92, 0.95], PR-AUC 0.876,
% Acc 0.883, F1 0.828. Eval-parity: per-ped 0.914 ~ per-window 0.913 (not an easier-eval
% artifact). Issue 4. Note convergence moved from a suspicious epoch 3 (leaky) to ~17.>

\subsection{Comparison with Published Baselines}
% <NOTE: insert Table~\ref{tab:baselines}. Ours = top AUC of the standard-protocol band
% with 2 streams; honest mid-pack Acc. Explain the high-AUC/mid-Acc profile (threshold-
% free ROC/PR-AUC are the comparable metrics; recall-favoring operating point). Issue 3.
% Both our rows (BiLSTM AND Transformer) lead the table on AUC with only 2 streams --
% the parsimony finding holds across two independent architecture families, not one.>
\begin{table}[H]
\caption{Crossing-prediction performance on PIE (standard protocol: train set01/02/04,
val set05/06, test set03). \label{tab:baselines}}
\begin{tabular}{lllccc}
\toprule
\textbf{Method} & \textbf{Venue/Year} & \textbf{Modalities (streams)} & \textbf{Acc} & \textbf{AUC} & \textbf{F1} \\
\midrule
PCPA \cite{kotseruba2021benchmark} & WACV 2021 & bbox + pose + context + speed (4) & 0.87 & 0.86 & 0.77 \\
Pedestrian Graph+ \cite{cadena2022pedgraph} & 2022 & pose graph + ego (2--3) & 0.89 & 0.90 & 0.81 \\
IntFormer % <NOTE: bib entry pending -- Issue 3 BibTeX pass>
 & 2021 & multimodal & 0.89 & 0.92 & 0.81 \\
PIT % <NOTE: bib entry pending -- Issue 3 BibTeX pass>
 & 2023 & multimodal & 0.91 & 0.92 & 0.82 \\
BiPed / PedFormer \cite{rasouli2023pedformer} & 2023 & multimodal & 0.91 & 0.90 & 0.85 \\
GTransPDM \cite{gao2024gtranspdm} & arXiv 2024 & bbox + pose + ego motion (3) & 0.90 & 0.87 & 0.82 \\
PIP-Net \cite{azarmi2025pipnet} & IEEE T-ITS 2025 & bbox, pose, speed, context, opt-flow, semseg, depth (7) & 0.915 & 0.897 & 0.846 \\
\midrule
BiLSTM (ours) & 2026 & \textbf{bbox + ego-speed (2)} & 0.883 & 0.932 & 0.828 \\
Transformer (ours, searched) & 2026 & \textbf{bbox + ego-speed (2)} & 0.894 & \textbf{0.950} & 0.845 \\
\bottomrule
\end{tabular}
% <NOTE: IntFormer and PIT rows still need bib keys (Issue 3's pending BibTeX/DOI
% pass -- arXiv ids for the field are in journal_prep/issue3_baseline_comparison/
% 03_baseline_comparison.md's Sources section). GTransPDM's own Table I lists these
% four via its own citation; consider citing GTransPDM's Table I as the secondary
% source for PIT/IntFormer/Pedestrian-Graph+/BiPed if their own papers remain
% unresolved by submission -- state this explicitly in the table note, don't hide it.
% Transformer row: staged 78-config search (val-only selection, test touched once),
% transformer/PLAN.md; 10k paired bootstrap vs the BiLSTM on the identical 2,094 test
% windows gives Delta-AUC = +0.0135, 95% CI [+0.0097, +0.0174] (excludes 0), paired
% t-test p=0.025 -- see transformer/phase5_analysis/05_comparison_report.md.>
\end{table}

\subsection{Feature Ablation: Ego-Speed Dominance}
% <NOTE: the headline ablation. baseline (bbox+ego) 0.932 vs bbox-only 0.753 (-0.18)
% vs attention 0.925 (no benefit). Ego-vehicle speed is the dominant predictor. Issue 2
% (05_variant_comparison.md). Figure or inline.>

\subsection{Cross-Set Generalization (LOSO)}
% <NOTE: leave-one-set-out 6-fold AUC 0.928 +/- 0.041; set03 fold 0.931 ~ fixed split
% -> set03 is representative, not an easy fold. Issue 5. Optional small table.>

\subsection{Observation Window and Prediction Horizon}
% <NOTE: window length insensitive (obs 8/16/30 -> 0.931/0.933/0.937, within seed
% noise); prediction horizon DOES matter (TTE 30/45/60 -> 0.960/0.948/0.919, every
% pairwise p<=0.008), CONFIRMED on a matched cohort (sample effect <=0.002). Issue 6.
% Insert Fig 06_ablation_figure.png (+ 06b matched).>

\subsection{Model Capacity}
% <NOTE: hidden 64/128/256 -> 0.927/0.933/0.938 (256 n.s. at 3.8x params); depth 1/2/3
% ~ 0.930/0.932/0.931. Model is small-data-limited; 128/2-layer justified. Issue 7 (+07b).
% Contrast worth drawing out in prose: this null result for the BiLSTM (more capacity
% doesn't help) is NOT universal across architectures -- Section~\ref{sec:transformer-
% comparison}'s Transformer search DID find that more capacity helps (4 layers beats 2,
% 794,241 vs the BiLSTM's 594,561 params) once the architecture family changes. Small
% data limits capacity gains for THIS recurrent architecture specifically, not for any
% sequence model over this input.>

\subsection{Transformer Comparison (Model-Choice Question, Answered Empirically)}
\label{sec:transformer-comparison}
% <NOTE: supervisor-requested extension (transformer/). Frames the "why not attention"
% question every reviewer of a small-data BiLSTM paper asks, and answers it with a
% measured result instead of a small-data argument alone.
%
% Protocol: identical frozen data/splits/pos_weight/seeds to the BiLSTM (transformer/
% PLAN.md Section 2); only the transformer's own architecture+recipe were searched --
% staged val-only selection (Stage A arch 36 configs -> Stage B recipe 36 -> transfer
% check 6 -> Stage C top candidates x5 seeds), test set03 touched exactly once
% (Stage D), mirroring Issue 8's own val-only discipline for the BiLSTM but at >=2x its
% search budget (78 vs 36 seed-42 configs).
%
% Headline (5 seeds, test set03 N=2094): BiLSTM 0.932 +/- 0.011 vs Transformer-searched
% 0.9497 +/- 0.0025 (794,241 params, d_model=128, num_layers=4, last-token pooling,
% sinusoidal positional encoding). 10k paired bootstrap of Delta-AUC on the identical
% windows (same resample indices both models, isolating sampling noise from genuine
% difference): Delta = +0.0135, 95% CI [+0.0097, +0.0174] -- excludes 0. Paired t-test
% over the 5 seeds: p=0.025 (n=5, stated as low-power; the window-paired bootstrap is
% primary evidence). BOTH pre-registered WIN conditions met -> Transformer adopted as
% co-headline model for this comparison, BiLSTM retained as the efficient baseline.
%
% The finding that makes this credible, not a coin-flip: Transformer-default (same
% architecture family, but trained with the BiLSTM's own un-searched recipe, zero
% tuning) is a statistical TIE with the BiLSTM (Delta=+0.0005, 95% CI [-0.0034,+0.0043]
% includes 0, p=0.83). Attention does not trivially beat recurrence on this 16x5
% signal -- a deliberately budgeted search is what found the improvement, echoing
% Issue 8's own point that a documented search should be trusted over an asserted
% architecture, in either direction (Issue 8: search CONFIRMED the BiLSTM's hand-set
% config; here it FOUND something better -- both outcomes are honest depending on what
% the search actually finds).
%
% LOSO (winner config, Issue-5 protocol, 6-fold): 0.939 +/- 0.044 vs the BiLSTM's
% 0.928 +/- 0.041 -- directionally consistent, though 6 folds is too few for a
% hypothesis test; reported as a generalization sanity check only.
%
% Latency (Issue 9 protocol, same M4 deployment hardware): the searched Transformer is
% FASTER than the BiLSTM per window at CPU batch-1 (0.457 ms vs 0.575 ms, ~1.3x more
% parameters notwithstanding) -- fully-parallel self-attention over T=16 tokens
% apparently outruns the BiLSTM's sequential recurrence on this hardware/batch size.
% Neither model is a latency concern for the live pipeline, which remains
% detection-bound (Section on Inference Latency, Issue 9).
%
% Full evidence: transformer/PLAN.md (pre-registered design + all DONE blocks),
% transformer/phase5_analysis/05_comparison_report.md (the verdict + parity gate),
% 04_final_summary.md, 06_latency_report.md, 07_loso_report.md. Insert
% transformer/phase5_analysis/05_comparison_figure.png as a figure here.>

\subsection{Inference Latency}
% <NOTE: isolated BiLSTM 0.575 ms/window (CPU), ~58x inside a 30 fps budget; CPU > MPS
% at batch 1 (GPU dispatch overhead). Pipeline detection-bound: YOLO26-M 33.7 ms (93%),
% BiLSTM 1.6 ms (4.5%) -> 27.5 fps. Issue 9. Insert Fig 09_latency_figure.png.>

\subsection{Detector-in-the-Loop Robustness}
% <NOTE: GT-box vs YOLO-box on demo clips (98 peds). AUC 0.962->0.953 (per-window),
% 0.958->0.948 (per-ped); decisions flip in 3%. Prediction robust to box noise. Weak
% links are perception: detector recall 88%, ByteTrack fragmentation (track purity 39%).
% Indicative (2 clips). Issue 10. Insert Fig 10_gt_vs_detector_figure.png.>

% =====================================================================
\section{Discussion}
% <NOTE: interpret. (1) Parsimony: 2 cheap streams reach top AUC -> heavy multimodal
% pipelines may be over-engineered for crossing prediction. (2) Why high-AUC/mid-Acc is
% honest, not "easy eval" (leakage removed, eval-parity, CIs, LOSO). (3) Positioning vs
% each prior work -- compress 04_positioning_vs_prior_work.md into a paragraph or table.
% (4) Deployment: the bottleneck is detection+tracking, not the predictor.>
\subsection{Limitations}
% <NOTE: state before a reviewer does. (a) Ego-speed carries much of the signal and
% partly encodes the ego-driver's anticipation (instrumented car slows for expected
% crossers) -- a legitimate inference-time signal but not purely vision-based; candidate
% speed-perturbation robustness check. (b) ByteTrack fragments identities -> needs re-ID
% in deployment. (c) Issue 10 is indicative (2 clips, 98 peds). (d) Single dataset (PIE).>

% =====================================================================
\section{Conclusions}
% <NOTE: restate the 3 contributions + the headline (0.932 at 0.575 ms with 2 streams,
% leakage-free, rigorously validated). Future work: re-ID-strengthened pipeline,
% speed-perturbation robustness, cross-dataset (JAAD), more demo clips for Issue 10.>

% =====================================================================
% BACK MATTER (use the template's macros)
% \authorcontributions{<NOTE>}
% \funding{<NOTE or "This research received no external funding.">}
% \dataavailability{The PIE dataset is publicly available \cite{rasouli2019pie};
%   our code and extraction protocol are available at <repo URL>.}
% \conflictsofinterest{The authors declare no conflict of interest.}
% \abbreviations{ <NOTE: AUC, BiLSTM, TTE, LOSO, ADAS, OBD, IoU ...> }
